The implementation follows the ones given by \citet{nori2024}.

For a set of particles $i$ with positions $\vec{x}_i$, masses $m_i$,
and smoothing length $h_i$, we compute the density for each particle as:
%
\begin{equation}
  \rho_i \equiv \rho(\vec{x}_i) = \sum_j m_j W(\vec{x}_{ij}, h_i).
  \label{eq:uldm:rho}
\end{equation}
and the derivative of its density with respect to $h$:

\noindent
Several possibilities exists to write the density gradient.
The standard way writes:
\begin{equation}
  \label{eq:uldm:grad_rho}
  \vec{\nabla} \rho_i =  \sum_j \frac{m_j}{\rho_j} \left( \rho_j - \rho_i \right) \vec{\nabla} W(\vec{x}_{ij}, h_i),
\end{equation}
which includes the substraction of the first order term. 
Another choice is \citep{nori2024}:
%
\begin{equation}
  \label{eq:uldm:grad_rho}
  \vec{\nabla} \rho_i =  \sum_j \frac{m_j}{\sqrt{\rho_i \rho_j}} \left( \rho_j - \rho_i \right) \vec{\nabla} W(\vec{x}_{ij}, h_i).
\end{equation}
This latter provides slightly better results on the 1D density front.

\noindent
As for the gradient density, several choices to write the Laplacian exits.
The simplest formulation  \citep[][Eq.~86]{price2012}:
%
\begin{equation}
  \nabla^2 \rho_i =  \sum_j \frac{m_j}{\rho_j} \rho_j \vec{\nabla}^2 W(\vec{x}_{ij}, h_i) =  \sum_j m_j  \vec{\nabla}^2 W(\vec{x}_{ij}, h_i),
\end{equation}
or removing the first order term \citep[][Eq.~90]{price2012}:
%
\begin{equation}
  \nabla^2 \rho_i =  \sum_j \frac{m_j}{\rho_j} \left( \rho_j - \rho_i \right)  \vec{\nabla}^2 W(\vec{x}_{ij}, h_i).
\end{equation}
Note that those two formulations totally fails the 1D density front test.
An alternative is to compute the Laplacian from the gradient (this works only in 1D I guess).
%
\begin{equation}
  \nabla^2 \rho_i =  \sum_j \frac{m_j}{\sqrt{\rho_i \rho_j}} \left( \vec{\nabla}\rho_j - \vec{\nabla}\rho_i \right) \vec{\nabla} W(\vec{x}_{ij}, h_i).
\end{equation}
This formulation is working.
Another formulation is \citep[][Eq.~18]{nori2018} {\bf (check with Price eq. 91)}:
%
\begin{equation}
  \nabla^2 \rho_i =  -2 \sum_j m_j \frac{ \rho_j - \rho_i }{\rho_j} \frac{ \vec{r_{ij}} \vec{\nabla} W(\vec{x}_{ij}, h_i)}{r^2_{ij}}.
\end{equation}





The Laplacian of the density write:
%
\begin{equation}
  \label{eq:uldm:laplacian_rho}
  \nabla^2 \rho_i =  \sum_j m_j \vec{\nabla}^2 W(\vec{x}_{ij}, h_i) \frac{\rho_j - \rho_i}{\sqrt{\rho_i \rho_j}} - \frac{|\vec{\nabla}\rho_i|^2}{\rho_i}.
\end{equation}

\noindent
The quantum potential of the particle $i$ is :
%
\begin{equation}
  \label{eq:uldm:grad_Q}
  Q_i =  \frac{\hbar}{2 m^2_\chi }  \sum_j \frac{m_j}{\rho_j} W(\vec{x}_{ij}, h_i)
  \left( \frac{\nabla^2\rho_j}{2\rho_j} -\frac{|\vec{\nabla}\rho_j|^2}{4\rho^2_j} \right)
\end{equation}


The acceleration of the particle $i$ is equal to the gradient of the quantum potential:

\begin{equation}
  \label{eq:uldm:grad_Q}
  \vec{\nabla{Q_i}} =  \frac{\hbar}{2 m^2_\chi }  \sum_j \frac{m_j}{\rho_j}\vec{\nabla} W(\vec{x}_{ij}, h_i)
  \left( \frac{\nabla^2\rho_j}{2\rho_j} -\frac{|\vec{\nabla}\rho_j|^2}{4\rho^2_j} \right)
\end{equation}



